\section{Introduction}\label{sec:introduction}

Campana's classification program separates compact K\"ahler geometry
into special varieties and bases of general type. Specialness records
the absence of differential forms that would detect such a base. One
of the program's central predictions is topological: the fundamental
group of a special compact K\"ahler manifold should become abelian
after passage to a finite unramified cover
\cite[Conjecture~A.1]{ClaudonHoring2013}.

We first specify the notion of specialness. For a holomorphic line
bundle \(L\) on a compact complex manifold \(X\), its Iitaka dimension
\(\kappa(X,L)\) is \(-\infty\) if no positive tensor power has a
nonzero section, and otherwise is the largest dimension of the images
of the meromorphic maps defined by its complete plurisection systems.
A manifold \(X\) of dimension \(n\) is \emph{special} if there is no
integer \(1\le p\le n\) and no nonzero sheaf morphism
\[
 L\longrightarrow\Omega_X^p
 \qquad\text{with}\qquad \kappa(X,L)=p.
\]
A line bundle with such a morphism and equality is a
\emph{Bogomolov sheaf}. A point is special. We prove the following
form of Campana's abelianity conjecture.

\begin{theorem}\label{thm:abelianity}
Let \(X\) be a connected smooth compact K\"ahler manifold of arbitrary
complex dimension. If \(X\) is special, then \(\pi_1(X)\) has an abelian
subgroup of finite index. Equivalently, \(X\) has a connected finite
unramified holomorphic cover with abelian fundamental group.
\end{theorem}

This is a positive resolution of the conjecture in the compact
K\"ahler setting. The conclusion concerns the entire discrete
topological fundamental group. It requires no linearity, residual
finiteness, projectivity, or prescribed Kodaira dimension. The finite
index may depend on \(X\), and torsion in \(\pi_1(X)\) is allowed.

\subsection{Kodaira dimension zero}

Campana's Theorem~6.7, applied to the zero-boundary orbifold
\((X|0)\), shows that a smooth connected compact K\"ahler manifold
of Kodaira dimension zero is special
\cite[Theorem~6.7, p.~92]{CampanaOrbifolds}. Combining this implication
with Theorem~\ref{thm:abelianity} gives the following consequence.

\begin{corollary}[Fundamental groups in Kodaira dimension zero]
\label{cor:kappa-zero-abelianity}
Let \(X\) be a smooth compact connected K\"ahler manifold with
\(\kappa(X)=\kappa(X,K_X)=0\). Then its ordinary topological fundamental
group \(\pi_1(X)\) is virtually abelian.
\end{corollary}

\begin{proof}
By Campana's established specialness implication, \(X\) is special.
Theorem~\ref{thm:abelianity} therefore applies.
\end{proof}

The specialness implication is also recorded in the orbifold core
formulation of \cite[Corollary \emph{Kodaira dimension along the core}]{OpenAIIitaka2026},
where its established provenance is noted. No projectivity,
vanishing of \(c_1(X)\), linearity, or residual-finiteness assumption
is required for Corollary~\ref{cor:kappa-zero-abelianity}.

\subsection{Conjecture S and affine-space uniformization}

We also record two geometric consequences identified by Campana.
For a connected compact K\"ahler manifold \(X\), let
\(\gamma_X:X\dashrightarrow\Gamma(X)\) be its ordinary
\(\Gamma\)-reduction. An irreducible compact analytic subspace \(Z\)
through a very general point is contained in the corresponding fibre
precisely when \(\pi_1(\widehat Z)\to\pi_1(X)\) has finite image,
where \(\widehat Z\) is the normalization of \(Z\). Set
\(\gamma d(X)=\dim\Gamma(X)\). This is the reduction for the identity
quotient of \(\pi_1(X)\), and maximal ordinary \(\Gamma\)-dimension
means \(\gamma d(X)=\dim X\)
\cite[Definition~2.2 and the paragraph preceding Conjecture~A.3]{ClaudonHoring2013}.
The two cases below are, respectively, Campana's Conjecture~S and the
bimeromorphic formulation of Iitaka's conjecture in his appendix
\cite[Conjecture~A.3 and Remark~A.4(2)]{ClaudonHoring2013}.

\begin{corollary}[Conjecture S and bimeromorphic Iitaka uniformization]
\label{cor:conjecture-s-iitaka}
Let \(X\) be a connected compact K\"ahler manifold of complex dimension
\(n\). Suppose that at least one of the following holds:
\begin{enumerate}[label=\textup{(\roman*)}]
\item \(X\) is special and \(\gamma d(X)=n\);
\item the ordinary universal cover of \(X\) is biholomorphic to \(\C^n\).
\end{enumerate}
Then \(X\) has a connected finite \'etale cover \(X'\to X\) such that
\(X'\) is bimeromorphic to a complex torus.
\end{corollary}

\begin{proof}
Theorem~\ref{thm:abelianity} supplies the Abelianity Conjecture.
Campana proves that the Abelianity Conjecture implies Conjecture~S in
\cite[Remark~A.4(3)]{ClaudonHoring2013}. We spell out the
finite-cover step. A finitely generated virtually abelian group has
a torsion-free abelian subgroup of finite index, so take the
corresponding connected finite \'etale cover \(X'\).
Specialness persists on \(X'\). Maximal ordinary
\(\Gamma\)-dimension persists as well. A positive-dimensional compact
subvariety through a very general point with finite group image
upstairs would have an image of the same dimension downstairs.
The induced map between the normalizations is finite \'etale, so
their fundamental-group images differ by finite index. The image
downstairs would therefore have finite group image, contrary to the
\(\Gamma\)-reduction criterion.

The Albanese map of \(X'\) is surjective with connected fibres by
Theorem~\ref{thm:specialness-inputs}. Since \(\pi_1(X')\) is
torsion-free abelian, its map to the Albanese lattice
\(H_1(X',\Z)/\mathrm{torsion}\) is an isomorphism. The group of a
smooth general Albanese fibre therefore has trivial image in
\(\pi_1(X')\). Maximal ordinary \(\Gamma\)-dimension forces that
fibre to have dimension zero. Connectedness makes the general fibre
one point, so the Albanese map is bimeromorphic to its torus target.
This proves~\textup{(i)}. Campana's second implication,
Conjecture~S to the stated bimeromorphic Iitaka conclusion, is
\cite[Remark~A.4(2)]{ClaudonHoring2013}; it gives~\textup{(ii)}.
\end{proof}

\subsection{Holomorphic convexity of the universal cover}

The compact K\"ahler linear Shafarevich theorem gives a further
geometric consequence. Recall that a complex manifold is
\emph{holomorphically convex} if the holomorphic hull of every nonempty
compact subset is compact.

\begin{corollary}[Holomorphic convexity for special manifolds]
\label{cor:special-shafarevich}
Let \(X\) be a connected smooth compact K\"ahler manifold. If \(X\) is
special, then its ordinary universal cover is holomorphically convex.
\end{corollary}

\begin{proof}
The group \(G=\pi_1(X)\) is finitely generated because \(X\) is a compact
smooth manifold. By Theorem~\ref{thm:abelianity}, it has an abelian
subgroup \(H\) of finite index, which is also finitely generated.
Write \(H\cong\Z^r\oplus T\) with \(T\) finite. Decomposing \(T\) into
cyclic factors gives a faithful finite-dimensional complex
representation \(\sigma:H\to\operatorname{GL}(V)\): represent each
infinite cyclic factor by powers of \(2\) on its own coordinate and
each finite cyclic factor by a primitive root of the corresponding
order on its own coordinate.
For the trivial group, use the one-dimensional trivial representation.

The induced representation \(\operatorname{Ind}_H^G\sigma\) is finite
dimensional because \([G:H]<\infty\), and it is faithful. Indeed, in
its direct-sum realization indexed by the left cosets \(G/H\), an
element of its kernel must preserve the summand indexed by \(H\), so
it belongs to \(H\); its action on that summand is given by the faithful
representation \(\sigma\). Thus \(G\) admits a faithful finite-dimensional
complex linear representation. The compact K\"ahler linear Shafarevich theorem
of Campana--Claudon--Eyssidieux
\cite[author's version, Corollary~5.9]{CCE2015} now gives the claimed
holomorphic convexity of the ordinary universal cover.
\end{proof}

No projectivity or maximal ordinary \(\Gamma\)-dimension is required
here. The conclusion is holomorphic convexity, not Steinness: for
example, \(\PP^1\) is special and its universal cover is compact.
In contrast, the smooth projective surface constructed in
\cite{OpenAIShafarevich2026} has a universal cover that is not
holomorphically convex. That surface is therefore not special.

\subsection{Compactifiable covers}

Claudon--H\"oring's compactifiable-cover results give another
consequence. Here a \emph{Zariski-open} embedding is holomorphic and
has closed analytic complement. The two parts require different
compactifications.

\begin{corollary}[Analytically compactifiable covers]
\label{cor:compactifiable-covers}
\begin{enumerate}[label=\textup{(\roman*)}]
\item Let \(X\) be a connected normal compact K\"ahler space and
\(p:V\to X\) a connected infinite \'etale Galois covering of complex
spaces, with deck group \(\Gamma\). If \(V\) embeds as a Zariski-open
subset of a normal compact complex space, then \(\Gamma\) is virtually
abelian.
\item Let \(X\) be a connected compact K\"ahler manifold and \(U\) its
ordinary universal cover. If \(U\) embeds as a Zariski-open subset of a
normal compact K\"ahler space, then there is a connected finite \'etale
cover \(X'\to X\) whose Albanese map
\[
 \alpha_{X'}:X'\longrightarrow A=\Alb(X')
\]
is a locally trivial holomorphic fibre bundle with simply connected
compact K\"ahler fibre \(F\). Writing \(q=\dim_{\C}A\), there is a
biholomorphism
\[
 U\simeq F\times\C^q.
\]
A point is allowed as \(A\) or \(F\).
\end{enumerate}
\end{corollary}

\begin{proof}
Suppose (i) has a counterexample and choose one of minimal base
dimension. The cover is normal because it is \'etale over normal
\(X\); normalizing an analytic compactification leaves this open set
unchanged. Claudon--H\"oring's minimal-counterexample reduction
\cite[Proposition~1.4]{ClaudonHoring2013} says that this minimal base is
smooth and special. Theorem~\ref{thm:abelianity} makes its fundamental
group virtually abelian. Its quotient \(\Gamma\) is then virtually
abelian, a contradiction.

For (ii), if \(\pi_1(X)\) is infinite, apply (i) to the ordinary
universal cover; if it is finite, it is already virtually abelian.
The K\"ahler compactification hypothesis is the one in
\cite[Definition~2.3]{ClaudonHoring2013}, so
\cite[Theorem~1.5]{ClaudonHoring2013} gives the asserted Albanese bundle
on a connected finite \'etale cover \(X'\). Pull this bundle back along
the universal cover \(\C^q\to A\). The automorphism group of the
compact K\"ahler fibre \(F\) is a complex Lie group, as recalled in
\cite[Section~2.C]{ClaudonHoring2013}. The associated principal
\(\operatorname{Aut}(F)\)-bundle is topologically trivial over the
contractible base \(\C^q\), and the Oka--Grauert principle makes it
holomorphically trivial over this Stein base. Thus the pullback is
\(F\times\C^q\), the product consequence also recorded in
\cite[Conjecture~1.1]{ClaudonHoring2013}. Since \(F\) is simply
connected, the pullback is the
ordinary universal cover of \(X'\), and hence also of \(X\).
\end{proof}

In (i), the compactifying space need not be K\"ahler and the deck action
need not extend to it. For a nonuniversal cover, only the deck group
\(\Gamma\), not the entire fundamental group of \(X\), is controlled.
The stronger K\"ahler compactification is used in (ii): the finite
cover \(X'\) need not be a product, and its fibre \(F\) need not be
projective.

\subsection{Prior work and the remaining group-theoretic obstacle}

Campana introduced specialness together with its orbifold framework
and the special/general-type decomposition in \cite{Campana2004}.
The same work established restrictions on solvable and linear
quotients of special fundamental groups. An important advance for the
full group was the proof of the abelianity conjecture for compact
K\"ahler threefolds by Campana--Claudon
\cite[Theorem~1.1]{CampanaClaudonThreefolds}, using the geometry of
orbifold surfaces. Their later work proves specialness of general
Albanese fibres in the smooth projective setting and reduces the
projective conjecture to a finiteness question when all finite covers
have irregularity zero
\cite[author's version, Theorem~2.4 and Proposition~3.2]{CampanaClaudon}.
These results explain the role of finite covers and Albanese fibre
groups in the argument below. We do not assume specialness of the
general Albanese fibre in the arbitrary compact K\"ahler case.

A classical predecessor is the case \(c_1(X)=0\) in real cohomology:
Beauville's decomposition gives a finite \'etale cover that is a
product of a complex torus and simply connected factors, and therefore
gives virtual abelianity in that case
\cite[Section~5, Theorem~2]{Beauville1983}.

The linear case has a separate geometric structure theory.
Building on work of Zuo and Eyssidieux, Campana--Claudon--Eyssidieux
describe linear Shafarevich reductions
by torus fibrations over bases of general type after finite cover and
modification \cite[author's version, Theorem~6.5]{CCE2015}.
Their semisimple theorem gives a general-type base for the torsion-free
representations used here
\cite[author's version, Theorems~1 and~6.3]{CCE2015}.
This supplies one branch of our proof. An infinite group, however,
need not have an infinite finite-dimensional complex linear image.
Controlling these remaining quotients is essential for a statement
about the full fundamental group.

The relation between fundamental groups and holomorphic tensors has
several relevant predecessors. Brunebarbe--Klingler--Totaro prove that
an infinite finite-dimensional linear image over any field forces a
nonzero symmetric differential on a compact K\"ahler manifold
\cite[Theorem~0.1]{BrunebarbeKlinglerTotaro2013}.
Brunebarbe--Campana prove that vanishing of all positive symmetric
powers of all exterior cotangent powers forces simple connectedness;
their argument develops the use of \(L^2\) forms and the
Poincar\'e--Atiyah--Gromov method
\cite[Theorem~1.1]{BrunebarbeCampana2016}.
Specialness permits many nonzero tensors, so their vanishing
hypothesis differs from ours. The task here is to obtain sufficient
\emph{positivity} from tensors with unitary Hilbert space coefficients.
Only after constructing a big ordinary line bundle, whose Iitaka
dimension equals the dimension of the base, and establishing projectivity do we invoke the orbifold cotangent positivity theorem
of Campana--P\u aun \cite[Theorem~7.11]{CampanaPaun2019}.

The spectral step also belongs to the study of \(L^2\) invariants of
infinite coverings. Atiyah's index theorem and regular-coefficient
bundle interpretation \cite[Theorem~3.8 and Section~6.1]{Atiyah1976},
and Dodziuk's reduced \(L^2\) de Rham--Hodge theory
\cite{Dodziuk1977}, provide foundational models for working on a
cover through Hilbert coefficients downstairs. Lott surveys the
zero-in-the-spectrum question and several geometric cases
\cite{Lott1996}. The unrestricted Riemannian assertion has
counterexamples \cite{FarberWeinberger2001}; the K\"ahler spectral
statement needed here is proved in Section~\ref{sec:spectral}.

A second issue is geometric descent over a possibly nonprojective
torus. Campana's correction to the Albanese theorem preserves
surjectivity and connectedness for special compact K\"ahler
manifolds, but its repaired assertion excluding multiple fibres
requires a projective target \cite[Theorem~1.2]{CampanaCorrection}.
Our proof therefore keeps the actual orders of meridians in the
chosen group quotient, including on exceptional divisors of new
models. Juanyong Wang's subadditivity and Albanese theorems provide
the required geometric input over arbitrary complex tori
\cite[Theorems~A(II) and~C]{JWang2021}. They apply to the smooth klt
pairs arising here: the boundary has simple normal crossings and
rational coefficients strictly less than one.

\subsection{Overview of the proof}

\paragraph{The group to be excluded.}
Finite \'etale covers preserve specialness. Since the
Albanese map of each cover is surjective, its first Betti number is
bounded by twice the dimension of \(X\). Pass to a finite cover on
which that number is maximal. Write
\(\alpha:X\to A=\Alb(X)\), and let \(N\) be the image in
\(G=\pi_1(X)\) of the fundamental group of a smooth general
Albanese fibre. We prove
\[
 1\longrightarrow N\longrightarrow G\longrightarrow
 \pi_1(A)\longrightarrow1,
\]
and show that \(N\) is finitely generated and \emph{FAb}: every
finite-index subgroup of \(N\) has finite abelianization.
The latter step uses Botong Wang's torsion theorem for isolated
characters in cohomology jump loci \cite[Theorem~1.3]{BWang2016}.
It remains to exclude infinite \(N\); a finite kernel of a lattice
quotient gives virtual abelianity by elementary group theory.

\paragraph{A minimal base carrying an infinite quotient.}
Assuming \(N\) infinite, consider fibrations of Albanese fibres on
all finite covers for which the quotient by the fibre-group image
is still infinite. Minimize the dimension \(\ell>0\) of their
bases. Compact analytic cycle spaces and Campana's meromorphic
\(\Gamma\)-reduction \cite{Barlet1975,CampanaGamma} globalize a
minimizer to a factorization
\[
 X\dashrightarrow Z\longrightarrow A.
\]
If \(B_0\) is the actual group image of a general fibre of the first
map, then \(Q=N/B_0\) is an infinite subgroup of \(R=G/B_0\), and
\[
 1\longrightarrow Q\longrightarrow R\longrightarrow\pi_1(A)
 \longrightarrow1.
\]
A very general fibre \(Y\) of \(Z\to A\) has dimension \(\ell\).
The actual meridian orders in \(R\) restrict to the same orders in
\(Q\) on \(Y\). For each boundary divisor \(D\subset Y\), denote
this finite order by \(m_D\), and put
\(\Delta_Y=\sum_D(1-m_D^{-1})D\).
The resulting orbifold \((Y,\Delta_Y)\) maps onto \(Q\) on
fundamental groups. Minimality holds simultaneously after the
compatible finite covers and for every further infinite quotient.
This is the input to the two branches in Figure~\ref{fig:proof}.

\paragraph{Positivity in the absence of an infinite linear image.}
Theorem~\ref{thm:zero-spectrum} first shows that an infinite normal
cover of a compact K\"ahler orbifold has zero in the spectrum of
its scalar Dolbeault Laplacian, with ordinary complex-valued
coefficients, in \emph{some} holomorphic degree.
The argument places the compact diagonal on
\((U\times U)/P\), where \(U\) is the cover and its deck
group \(P\) acts diagonally. The diagonal \(\{(u,u):u\in U\}\)
therefore has compact quotient. Under a contrary spectral gap, weighted solutions and
finite-propagation smoothing preserve the nonzero cohomology class
of the diagonal after push-down, while making its pairing tend to
zero. The contradiction produces the required low spectrum.

When every finite-dimensional complex linear image of \(Q\) is
finite, we construct nonzero holomorphic tensors with unitary Hilbert
coefficients. We use the spectral construction in the nonamenable
case and Shalom's reduced-cohomology theorem in the amenable case
\cite[Theorem~4.2]{Shalom2006}. Taking exterior powers gives tensors with rank-one coefficient
image; these lines define maps to Hilbert projective space. We minimize the rank among such
maps, prove discreteness of the monodromy on their regular image
germs, and compactify the fibres using transverse volume bounds.
A rank smaller than \(\ell\) would contradict the minimal-base
condition. Full rank yields a big cotangent line, and hence
\(K_Y+\Delta_Y\) is big by Theorem~\ref{thm:hilbert-positivity}.
When \(Q\) instead has an infinite linear image, its FAb property
produces the required semisimple image and minimality gives generic
largeness. The general-type criterion of
Campana--Claudon--Eyssidieux then gives \(K_Y\) big.

\paragraph{Descent and the contradiction.}
Both branches give fibre positivity on every sufficiently prepared
model. Theorem~\ref{thm:torus-descent} turns it into a
positive-dimensional orbifold base of general type whose orbifold
pluriforms pull back
to a Bogomolov sheaf on a finite cover of \(X\). The key point is
to retain actual quotient meridian orders on each model; no
birational invariance of the associated log Kodaira dimension is
assumed. Differential vanishing cancels the boundary poles on a
neat model. Specialness excludes the resulting sheaf, so \(N\)
is finite and Theorem~\ref{thm:abelianity} follows.

\begin{figure}[htbp]
\centering
\begingroup
\small
\setlength{\fboxsep}{7pt}
\fbox{\begin{minipage}{0.9\linewidth}\centering
Infinite FAb group \(N\): minimize a base dimension \(\ell>0\).\\
Obtain \(Q=N/B_0\) and \((Y,\Delta_Y)\) over the Albanese torus.
\end{minipage}}
\par\smallskip
\(\Downarrow\)
\par\smallskip
\begin{tabular}{@{}c@{\hspace{0.025\linewidth}}c@{}}
\fbox{\begin{minipage}[t]{0.405\linewidth}\centering
\(Q\) has an infinite complex\\linear image in finite dimension.\par\smallskip
FAb, minimality, and the\\semisimple general-type criterion\par\smallskip
\(\Longrightarrow K_Y\) big
\end{minipage}}
&
\fbox{\begin{minipage}[t]{0.405\linewidth}\centering
Every finite-dimensional complex\\linear image of \(Q\) is finite.\par\smallskip
Hilbert tensors, discrete monodromy,\\and minimality\par\smallskip
\(\Longrightarrow K_Y+\Delta_Y\) big
\end{minipage}}
\end{tabular}
\par\smallskip
\(\Downarrow\)
\par\smallskip
\fbox{\begin{minipage}{0.9\linewidth}\centering
Torus descent with actual meridian orders\par\smallskip
A Bogomolov sheaf on a finite cover of \(X\): contradiction.
\end{minipage}}
\endgroup

\caption{The contradiction when the Albanese fibre-group image
\(N\) is infinite. The minimal quotient \(Q\) and the fibre
\((Y,\Delta_Y)\) are constructed together. Each branch must give
positivity on all sufficiently prepared models before torus descent
applies. The arrows show logical implications.}
\label{fig:proof}
\end{figure}

\paragraph{Organization.}
Section~\ref{sec:preliminaries} fixes specialness, group reductions,
and the precise external geometric inputs.
Section~\ref{sec:spectral} proves the scalar spectral theorem.
Section~\ref{sec:hilbert} constructs the Hilbert tensors and proves
their positivity theorem under the minimality condition.
Section~\ref{sec:descent} develops meridian bookkeeping and torus
descent. Section~\ref{sec:completion} establishes the Albanese exact
sequence and FAb property, constructs the minimal relative base,
and assembles the two representation cases to finish the proof.
